% Общая преамбула для решений и шпаргалки
\usepackage[T2A]{fontenc}
\usepackage[utf8]{inputenc}
\usepackage[russian]{babel}
\usepackage{amsmath,amssymb}
\usepackage{xcolor}
\usepackage[most]{tcolorbox}
\usepackage{enumitem}
\usepackage{needspace}

\setlist{nosep,leftmargin=*}
\setlength{\parindent}{0pt}
\setlength{\parskip}{3pt}

% --- обозначения ---
\newcommand{\dd}{\mathrm{d}}
\newcommand{\avg}[1]{\overline{#1}}
\newcommand{\vp}{v_{\text{в}}}            % наивероятнейшая скорость
\newcommand{\vav}{\bar v}                  % средняя скорость
\newcommand{\erf}{\operatorname{erf}}

% --- задача ---
% \begin{problem}{1.3}{dz}{условие} ... \end{problem}
% второй аргумент: dz — задача из ДЗ, ex — разобранный пример книги,
% exdz — пример, который одновременно входит в ДЗ
\definecolor{dzcol}{RGB}{180,40,40}
\definecolor{excol}{RGB}{40,90,170}
\makeatletter
\newcommand{\tag@dz}{\textcolor{dzcol}{ДЗ}}
\newcommand{\tag@ex}{\textcolor{excol}{пример}}
\newcommand{\tag@exdz}{\textcolor{excol}{пример}\,+\,\textcolor{dzcol}{ДЗ}}
\newcommand{\probtag}[1]{\csname tag@#1\endcsname}
\makeatother
\newtcolorbox{problem}[3]{enhanced,breakable,
  colback=white,colframe=black!55,boxrule=0.5pt,arc=1pt,
  left=4pt,right=4pt,top=3pt,bottom=3pt,before skip=6pt,after skip=6pt,
  fonttitle=\bfseries,coltitle=black,colbacktitle=black!8,
  title={Задача #1\hfill{\normalfont\small\bfseries\probtag{#2}}},
  before upper={\textit{#3}\par\smallskip}}

% блок опорных формул в начале параграфа
\newtcolorbox{basics}{enhanced,breakable,colback=blue!3,colframe=excol!70,
  boxrule=0.5pt,arc=1pt,left=4pt,right=4pt,top=3pt,bottom=3pt,
  fonttitle=\bfseries,title={Опорные формулы}}
\newcommand{\sub}[1]{\par\smallskip\textbf{#1}\par}

\newcommand{\idea}[1]{\textbf{Идея:} #1\par}
\newtcolorbox{answerbox}{colback=yellow!12,colframe=black!40,boxrule=0.4pt,arc=1pt,
  left=3pt,right=3pt,top=1pt,bottom=1pt,before skip=4pt,after skip=2pt}
\newcommand{\answer}[1]{\begin{answerbox}\textbf{Ответ:} #1\end{answerbox}}
\newcommand{\note}[1]{\par{\small\textcolor{black!60}{\textit{Замечание:} #1}}\par}
